\\newpage
\\section{误差分析}

\\subsection{预测模型误差}
\\begin{enumerate}
    \\item[(1)] \\textbf{图结构简化}：固定邻接矩阵无法反映交通流的动态传播特性，实际路网存在时变拓扑（如施工封闭）。
    \\item[(2)] \\textbf{特征工程}：手工设计的时空特征可能遗漏重要信息，端到端学习可以更充分地利用数据。
\\end{enumerate}
MAE在正常交通条件下为12.3辆/5min，在拥堵条件下上升至28.7辆/5min，模型对拥堵状态的预测误差较大。

\\subsection{控制策略误差}
\\begin{enumerate}
    \\item[(1)] \\textbf{模型简化}：信号控制模型忽略行人过街、公交优先等复杂因素，高峰时段延误估算偏低。
    \\item[(2)] \\textbf{延迟影响}：实际控制系统存在通信和计算延迟（约200-500ms），影响控制效果。
\\end{enumerate}
测试表明，在10\\%流量扰动下，协同控制策略仍能将总延误控制在基准值的1.15倍以内。
